\section{Introduction}\label{sec:intro}

\subsection{Context and problem}
A reconfigurable intelligent surface (RIS) shapes the radio channel with many
passive, electronically switched elements~\cite{basar2019,wu2019}. Most RIS work
keeps the configuration fixed during an OFDM symbol and changes it, if at all,
between symbols. A \emph{time-modulated} RIS (TM-RIS), also called a
space-time-coding metasurface, instead switches each element several times within
the symbol~\cite{zhang2018}. The reflected signal then contains harmonics of the
switching frequency, and the surface can steer, scramble or split energy across
frequency as well as space. OFDM-aligned time modulation has been proposed for
waveform security and directional modulation~\cite{xu2023tmirs,tao2025gflownet,wang2026generative}
and for joint sensing and security~\cite{tao2025isac}.

For the user that the surface serves, the same harmonics are a problem. A harmonic
at $h$ times the subcarrier spacing moves energy from transmitted subcarrier $m$ to
received subcarrier $m+h$, so the frequency-domain channel matrix is no longer
diagonal. Iudice \emph{et al.}~\cite{iudice2026} recently derived a physically
consistent model of this channel in which each element has a frequency-dispersive
(finite-memory) response, as a real varactor cell does. They showed that a
generalized cyclic prefix removes inter-symbol interference but that the coupling
between subcarriers remains for any prefix length, and that with dispersive elements
the coupling matrix is not even Toeplitz. Receiver design was explicitly left open.

How large is the problem? For a two-state element with ideal reflection $\pm1$ and
duty cycle $D$ (fraction of the period in state 1), Parseval's theorem gives the
fraction of reflected energy that stays on the original subcarrier as
$(2D-1)^2$ (Appendix~\ref{app:parseval}). At $D=0.75$ only 25\% stays; the
remaining 75\% is spread over the other harmonics. A conventional one-tap OFDM
receiver treats that 75\% as interference, so its signal-to-interference ratio is
at most $(2D-1)^2/(1-(2D-1)^2)=-4.8$\,dB, and its rate saturates however high the
SNR. In our evaluation the one-tap receiver delivers 0.42\,bit/s/Hz at both 20 and
30\,dB, against 6.6 and 9.9\,bit/s/Hz with perfect channel knowledge and joint
equalization (Section~\ref{sec:res-est}).

Time modulation does not, by itself, give the user more rate: for a single user, an optimized \emph{static} RIS
delivers more (Section~\ref{sec:res-motiv}). A TM-RIS is deployed for its other functions, such as directional
security, harmonic steering of several users or sensing, and the question this paper answers is how much of the
resulting loss the legitimate receiver can win back. Recovering this rate requires the full coupling matrix at the receiver. The matrix
has up to $M^2$ entries ($M=3276$ subcarriers in a 100\,MHz 5G NR carrier) and
depends on the propagation on both hops, on the switching pattern, and on the
element response. A generic estimator of a non-diagonal OFDM channel, such as the
basis-expansion model used for high-mobility channels~\cite{tang2007}, does not
know where the coupling comes from. It must estimate every harmonic separately.
Because rectangular switching leaves a slowly decaying $1/h^2$ energy tail, it
either truncates (and then misses 3.4\% of the energy at $|h|\le 12$, which caps its
accuracy near $-15$\,dB) or keeps all harmonics (and then needs many pilots).

\subsection{Central idea}
The receiver of a TM-RIS link is in an unusual position: it knows \emph{how} the
channel varies inside the symbol, because the variation is the switching pattern,
and the switching pattern is set by the network. What it does not know is the
propagation and the exact element response. This paper turns that observation
into an exact structural statement and builds the entire receiver on it.

Grouping the harmonic expansion of~\cite{iudice2026} by element \emph{state} rather
than by harmonic gives (Proposition~\ref{prop:sep})
\begin{equation*}
\bm H=\sum_{q=1}^{Q}\sum_{s=1}^{2}\underbrace{\bm T_{q,s}}_{\text{timing only}}\circ
\big(\underbrace{\bm a^{\rm d}_q}_{\text{user hop}}\,
(\underbrace{\bm\Gamma_s\odot\bm a^{\rm s}_q}_{\text{BS hop and dispersion}})^{\!\top}\big),
\end{equation*}
where each $\bm T_{q,s}$ is a Toeplitz matrix fixed by the switching timing alone.
The algebra is a regrouping of~\cite[eq.~(5)]{iudice2026}, and we credit it as such.
What is new is what the regrouping implies for a receiver. Each consequence becomes
one block of the receiver:
\begin{enumerate}
\item The element dispersion multiplies only the input side, together with the
unknown base-station (BS)--RIS channel, so it can be absorbed into an unknown
per-state input response. \emph{The receiver needs no element model}
(Section~\ref{sec:est}). Assuming the ideal $\pm1$ element of earlier work costs
1.4/6.6\,dB at 20/30\,dB; the timing-only receiver loses 0.08/0.11\,dB.
\item All elements see the same propagation paths, so the unknown factors live in a
low-dimensional delay subspace shared by all elements and both states. The unknowns
fall from up to $M^2$ to about $3Lr$ ($L$ delay functions, rank $r$), and two pilot
symbols suffice for a sparse channel (four to eight for 3GPP TDL-C at 100\,ns).
\item A clock offset between the RIS and the receiver changes $\bm T_{q,s}$ to
$\bm D\bm T_{q,s}\bm D^{*}$, which is \emph{exactly} a delay of the user hop and an
advance of the BS hop. This explains why a pilot-residual timing search fails at high
SNR (the delays absorb the offset), and why the Bayesian evidence, which also asks how
probable the absorbing delays are, recovers it (Section~\ref{sec:timing}).
\item The BS--RIS side does not move, so it can be learned once over many frames, and
each new frame estimates only the user side (Section~\ref{sec:twoscale}).
\item Each $\bm T_{q,s}\bm x$ is a linear convolution, so products with $\bm H$ cost
$\mathcal O(QM\log M)$ and the MMSE equalizer can run matrix-free
(Section~\ref{sec:eq}).
\item The $1/h$ amplitude envelope of $\bm T_{q,s}$ explains, quantitatively, why
truncated estimators and banded equalizers fail on this channel
(Corollary~\ref{cor:tail}).
\end{enumerate}

\subsection{Contributions}
\begin{itemize}
\item \textbf{C1, receiver structure.} To our knowledge, the first pilot-based receiver for OFDM through a
(dispersive) TM-RIS: pilot-based estimation of the complete coupling matrix and joint
equalization, organized around the timing/propagation separation above
(Fig.~\ref{fig:pipeline}, Algorithms~\ref{alg:est}--\ref{alg:cg}).
\item \textbf{C2, timing-only estimation.} Absorbing the dispersion into the unknown
input response, so that no element model, calibration or resonance knowledge is needed,
and a low-rank shared-subspace estimator that covers identical, cyclically shifted and
arbitrary per-element switching patterns in one form.
\item \textbf{C3, blind clock-offset recovery.} The opposite-delay identity for a
TM-RIS clock offset and its physical meaning (bulk hop delays cancel if the RIS follows the incident signal;
the BS--RIS delay becomes an offset if it follows the BS clock), the resulting overfitting mechanism of
residual-based timing, and an evidence-based estimator of the offset and the delay reference that absorbs the offset from two pilot symbols
with no rate loss against known timing. (Because of the identity, the offset is defined only
relative to the reference of the delay prior; the meaningful metric is the rate loss.)
\item \textbf{C4, exhaustive evaluation.} An evaluation against the strongest
baselines we could build (full-harmonic basis expansion given 16$\times$ more pilots,
sparse recovery, unstructured least squares, literature equalizers given the true
channel), on physical and 3GPP channels, with coding, mobility and pilot overhead,
hardware impairments, timing errors, multiple panels and interference, and the
complexity in operations and measured time. Pass/fail criteria were fixed before each
experiment; failures are reported with their diagnosis.
\end{itemize}

We adopt, and do not claim: low-rank modelling of RIS cascaded
channels~\cite{wang2020compressed}; the DPSS delay basis~\cite{slepian1978,zemen2005};
matrix-free iterative equalization through product-convolution
operators~\cite{hrycak2010} (here the representation is exact rather than a
basis-expansion approximation); learning the quasi-static BS--RIS channel on a slow
timescale~\cite{hu2019twotimescale}; Bayesian evidence for model and timing
selection~\cite{mackay1992,access2019bayes}; and the observation, known from
space-time-coding metasurfaces, that a time shift of the switching pattern multiplies
harmonic $h$ by a linear phase~\cite{zhang2018}. Section~\ref{sec:related} explains
how each differs from its use here.

\subsection{Organization and notation}
Section~\ref{sec:related} positions the work. Section~\ref{sec:model} gives the system
model and every assumption with its justification. Section~\ref{sec:structure} derives
the structure of the coupling. Section~\ref{sec:receiver} presents the receiver,
algorithm by algorithm, with its complexity. Section~\ref{sec:method} describes the
evaluation protocol and baselines, Section~\ref{sec:results} the results, and
Section~\ref{sec:tradeoffs} the trade-offs and limitations.

Bold lower- and upper-case letters denote vectors and matrices; $\circ$ and $\odot$ are
the Hadamard product of matrices and vectors; $(\cdot)^{\top}$, $(\cdot)^{\mathsf H}$ and
$(\cdot)^{*}$ are transpose, conjugate transpose and conjugate; $\diag(\bm a)$ is the
diagonal matrix with $\bm a$ on its diagonal; $\vect(\cdot)$ stacks columns;
$\mathcal{CN}(\bm\mu,\bm C)$ is the circularly symmetric complex Gaussian distribution;
$\mathrm{sinc}(x)=\sin(\pi x)/(\pi x)$. Confidence intervals are $\pm1.96\,s/\sqrt{n}$
over $n$ independent channel draws.
